\subsection{Design and primary contrasts}
\label{sec:f2-admission}
A separate confirmatory BIRD study crosses six independent builders with
three seeds under strategy forcing (a named mechanism from a fixed set
versus free invention) and conformance gating, giving four arms:
II-A (free, ungated), II-B (free, gated), II-C (forced, ungated), and
II-D (forced, gated). The evaluation covers
$1{,}169$ questions from nine databases never opened during
development (one disclosed $18$-item smoke-test exception), with the
official BIRD execution scorer and a pre-outcome-frozen analysis later
re-run as a versioned corrected reanalysis. The implemented gate bundles
nontrivial-mechanism structural screening with conformance verification
--- two of eight forced prompt-transformation slots cannot pass even
when neutrally valid --- so the treatment does not isolate verification
causally.

\begin{table}[H]
\centering\small
\begin{tabular}{@{}lrr@{}}
\toprule
Metric & II-A & II-D \\
\midrule
Candidate count & 5.78 & 3.83 \\
Oracle accuracy (\%) & 76.14 & 75.70 \\
Best-fixed accuracy (\%) & 67.58 & 67.30 \\
Headroom (pp) & 8.56 & 8.40 \\
\midrule
D$-$A headroom (pp) & \multicolumn{2}{c}{$\RevisionPrimaryDelta$} \\
Unadjusted bootstrap 95\% CI & \multicolumn{2}{c}{$\RevisionPrimaryCI$} \\
One-sided sign-flip $p$ & \multicolumn{2}{c}{$\RevisionPrimaryP$} \\
\bottomrule
\end{tabular}

\caption{\textbf{Corrected primary analysis.} Means over $18$
builder--seed cells on $1{,}169$ questions; \bare{} is included in both
the oracle and best-fixed comparators.}
\label{tab:primary}
\end{table}

The diagonal forced/gated minus free/ungated contrast in single-call-inclusive
oracle headroom is $\RevisionPrimaryDelta$\,\%, with cell contrasts
ranging over $\RevisionPrimaryRange$\,\%: the bundled screening policy
does not improve mean headroom. The original uncertainty analysis is
retained in Table~\ref{tab:primary}; this mean comparison does not
establish equivalence. Structural admission is therefore not evidence
that a mechanism activates on relevant tasks or adds stable utility.

\begin{figure}[H]
\centering
\includegraphics[width=0.72\linewidth]{figures/fig_results_revision.pdf}
\caption{\textbf{Versioned corrected reanalysis.} Left: full-set
D$-$A cell contrasts with the pooled 95\% interval. Right: four-arm
headroom on the shared core-$400$ tasks with a common \bare{} baseline.}
\label{fig:results}
\end{figure}

On the shared core-$400$ tasks the gate contrast is
$\RevisionGateDelta$\,\%, and the strategy and interaction secondaries
are likewise inconclusive (Table~\ref{tab:factorial}); the fixed-$K$
mean contrast $\RevisionKmatchedDelta$\,\% provides a population-size
sensitivity check. The open-mechanism arm ran only on
the development side and has no confirmatory estimand. These are all
single-run contrasts, and the execution variability reported in
the R2/R3 audits in Appendix~\ref{app:r2} ($8.54\%$ R2 verdict flips; $14.5\%$ of tasks
changing \bare{} correctness across three runs; a $6.50$\,\% same-code
one-run gap) characterizes them as coverage opportunities rather than
capability differences.

\begin{table}[H]
\centering\small
\begin{tabular}{@{}lrrrr@{}}
\toprule
Contrast & Effect (pp) & 95\% CI (pp) & Raw $p$ & Holm $p$ \\
\midrule
Gate & $-0.69$ & $[-1.56, +0.20]$ & 0.026 & 0.079 \\
Strategy & $+0.42$ & $[-0.32, +1.51]$ & 0.263 & 0.428 \\
Interaction & $-0.87$ & $[-2.83, +0.65]$ & 0.214 & 0.428 \\
\bottomrule
\end{tabular}

\caption{\textbf{Corrected core-$400$ factorial reanalysis.} Gate is
$[(B{-}A)+(D{-}C)]/2$, strategy $[(C{-}A)+(D{-}B)]/2$, interaction
$(D{-}C)-(B{-}A)$; raw and Holm-adjusted $p$-values.}
\label{tab:factorial}
\end{table}
